\section{Conclusion}\label{sec:conclusion}

The selected catalogue contains results proved within specified frameworks and physical proposals with different degrees of support. The exact computations and analytic tests do not justify a count of ``fatal errors'' in an author's work. The toroidal negative full indices and their subsequent hair interpretation are already in the literature. The caustic calculation concerns the classical dust approximation and leaves the homogeneous rolling solution intact. The divisor singularity belongs to the orientifold cover, while a local $I_0^*$ family shows why non-Higgsability alone does not prohibit weak coupling. Two of Sen's statements now hold in general form: his weak coupling family is the normal form of every degeneration with $v(j)<0$, and the entropy of a two-charge small black hole is fixed by one critical value, so that the Gauss--Bonnet truncation fails exactly from $D=6$ on and Prester's Lovelock completion is forced.

The remaining questions are mathematical and physical construction problems: the identification of the limit of Theorem~\ref{thm:normalform} with a type~IIB orientifold beyond its geometry, the near-horizon value $R^*$ of Theorem~\ref{thm:twocharge} for the complete heterotic effective action in $D\ge6$, the full monopole $L^2$ statement beyond the cited cases, general hair identification, a controlled definition of the quantum entropy function, and extensions of the stated soft and string-field results. They are not resolved by the finite tables. A meaningful new criticism would require a precise original statement with all its hypotheses retained and a counterexample within those hypotheses. No new false proof by Sen meeting that criterion is established in this revision.

A reader who wants to repeat the computations of this paper will find them in the archive \cite{Bha26}, together with the Lean formalisation of the finite statements: the Fourier coefficients of Tables~\ref{tab:phi10} and~\ref{tab:tor2}, the wall-crossing identities of Theorem~\ref{thm:wallcrossing} and the twenty-one negative indices of Theorem~\ref{thm:positivityfails}, each computed inside the Lean kernel from the definitions of the theta functions. The catalogue and physical interpretations are not formalised by these certificates. Mathematical parts of such subjects can be formalised, but that work must not be inferred from a finite arithmetic check.
